\section{Ograniczenia środowiska}
\label{sec:ograniczenia-srodowiska}

Choć zaprojektowany i~zaimplementowany model środowiska
Ultimate Texas Hold’em jest wystarczający do przeprowadzenia badań nad strategiami podejmowania decyzji, to nie odwzorowuje on wszystkich aspektów gry w~kasynie.

Prawdą jest, że zakłady Trips i~Progressive nie wprowadzają dodatkowego momentu decyzyjnego,
ponieważ ich wynik nie zależy od decyzji związanej z~akcją Play. Niemniej zakłady poboczne wpływają na
rozkład końcowych wypłat i~utrudniałyby interpretację jakości strategii,
a~Progressive wymagałby dodatkowo modelu puli utrzymywanej pomiędzy
rozdaniami, co naruszałoby niezależność epizodów. Oba zakłady poboczne można zaimplementować jako rozszerzenie warstwy rozliczeń, co może stanowić przedmiot dalszych badań.

Silnik nie modeluje również
salda gracza, sesji, ryzyka bankructwa ani zarządzania kapitałem, a~wartości
zakładów wyraża w~jednostkach Ante, dzięki czemu strategie można porównywać
niezależnie od nominału.

Środowisko zaprojektowano tak, że reprezentuje ono pojedynczego gracza wobec
krupiera i~nie symuluje pozostałych miejsc przy stole, a~zakryte karty innych
graczy i~tak byłyby nieznane, więc rozkład kart widocznych dla gracza
odpowiadałby losowaniu bez zwracania ze standardowej talii.

Obserwacja jest obiektem domenowym zawierającym karty, fazę i~zbiór
legalnych akcji. W~tej postaci jest czytelna i~wygodna w~testach, lecz nie może
być bezpośrednio podana modelom uczenia maszynowego. Wynika to z~tego, że kodowanie
\texttt{UTHObservation} do reprezentacji numerycznej pozostaje poza zakresem
klasy \texttt{UTHGame}, ponieważ porównanie reprezentacji stanu jest jednym
z~elementów badawczych pracy. Wspólny silnik jest wygodny, bo pozwala zestawiać różne
reprezentacje na identycznych rozdaniach i~regułach.

\section{Podsumowanie}
\label{sec:podsumowanie-srodowiska-uth}

W~niniejszym rozdziale przedstawiono projekt oraz implementację pojedynczego
rozdania Ultimate Texas Hold’em. Silnik podzielono na niezależne moduły
odpowiedzialne za reprezentację kart, ocenę układów, rozdawanie,
przechowywanie stanu, walidację akcji, porównanie układów oraz rozliczanie zakładów,
a~przebieg rozdania odwzorowano jako maszynę stanów obejmującą fazy
\texttt{PREFLOP}, \texttt{FLOP}, \texttt{RIVER} i~\texttt{TERMINAL}, w~której
każda akcja \texttt{BET} kończy część decyzyjną i~wymusza co najwyżej jeden
zakład Play.

Ważnym elementem projektu jest rozdzielenie pełnego stanu
\texttt{RoundState} od obserwacji \texttt{UTHObservation}, mianowicie agent otrzymuje
wyłącznie własne karty, odkryte karty wspólne, fazę i~legalne akcje, podczas
gdy karty krupiera, karty spalone i~kolejność talii pozostają wewnętrzne.
Porównanie układów korzysta z~ogólnego ewaluatora wybierającego najlepsze pięć kart
z~siedmiu, a~rozliczenia Ante, Blind i~Play są przechowywane dokładnie za
pomocą typu \texttt{Fraction}. Deterministyczne źródło \texttt{FixedDeck}
oraz opcjonalny mechanizm historii zapewniają powtarzalność i~możliwość
analizy przebiegu rozdań.

Na tym etapie powstał silnik gry, który zapewnia podstawę do implementacji agentów bazowych, infrastruktury eksperymentalnej oraz metod uczenia przez wzmacnianie analizowanych w~dalszej części pracy.
